\documentclass[11pt]{article}
\usepackage[T1]{fontenc}
\usepackage{lmodern,microtype}
\usepackage{amsmath,amssymb}
\usepackage[letterpaper,margin=1in]{geometry}
\usepackage{booktabs,array,enumitem,longtable,xcolor}
\usepackage[colorlinks=true,linkcolor=blue!55!black,urlcolor=blue!55!black]{hyperref}
\title{Revision log for Part IV, Draft 3\\
\large Response to Claude's referee report on Draft 2}
\author{Prepared by ChatGPT (OpenAI) for David Joyner\\
\normalsize Referee: Claude (Anthropic)}
\date{November 7, 2026 (Eastern Standard Time)}
\begin{document}
\maketitle
\section*{Source and scope}
This log records changes from \path{mjl_part4_draft2.tex}
(formatted title-page version) to \path{mjl_part4_draft3.tex}.
It responds to the attached \path{mjl_part4_draft2_referee_report_claude.tex}.
The arithmetic, representation-theoretic, and trace-formula claims
that have not been established remain explicitly conditional or conjectural.
The numerical references from Part III were checked against its compiled
\emph{Draft 9} rather than Draft 7: Lemma 3.3, Theorems 4.3, 7.6,
and 9.1, Corollaries 5.1, 5.2, and 7.8, and Proposition 11.3
retain their numbering.  Parts I and II were consulted in their
published-to-GitHub Draft 7 source versions when this edit began.
This is not an independent referee verification of every theorem in the series.

\section*{Disposition of substantive referee findings}
\begin{enumerate}[leftmargin=*,itemsep=0.55em]
\item \textbf{Referee \S1.1: Odd-$n$ star-map correction $u(x)=1$ needs justification.} Explain that $u(x)\in\{\pm1\}\cap\mu_n=\{1\}$, not merely that Zou's multiplicativity identity holds.  State scalar and automorphic central-character consequences.
\item \textbf{Referee \S1.2: Split regular element may have singular power; field regularity deserves proof.} In Remark 3.3 explain that $\mathrm{diag}(1,\zeta_3)$ is not good over $\mathbb Q_3(\zeta_3)$, and give the coprime-degree/Kummer proof for field tori.
\item \textbf{Referee \S1.3: General pushforward density formula in Proposition 4.2 is false.} Replace the proof by a compact-open principal-unit calculation. The result remains $n^{k-1}|n|_F^{-(r-1)}$. Include the counterexample $\mathbb Z\xrightarrow{\times n}\mathbb Z$ in Remark 4.3.
\item \textbf{Referee \S1.4: Name the split-place cover isomorphism.} Cite Part I, Proposition 2.6 explicitly in Theorem 5.2's proof.
\item \textbf{Referee \S1.5: Explain Zou's central-character hypothesis and remove redundant regularity proviso.} Proposition 7.1 states automatic regularity for a degree-$r$ field torus.  It explains the root sum when $\omega_\pi|_{\mu_n}$ is or is not trivial, and why an actual lift pair forces triviality.
\item \textbf{Referee \S1.6: Geometric Proposition 8.1 was nearly tautological.} Add the FK-compatible function class (Definition 8.1), prove its necessity (Lemma 8.2), and replace the global criterion with explicit split-place pointwise identities $J_{\widetilde H_v}(x^*,\widetilde f_v)=\lambda_v J_{H_v}(x,f_v)$ and the product condition $\prod_{v\notin S}\lambda_v\prod_{v\in S}b_{D,v}=1$.  State the unproved existence and singular-germ obligations.
\item \textbf{Referee \S1.7: Real-place statement about Li is misleading.} State that $K$ is totally imaginary and has no real archimedean places.  Note the splitting at complex places and cite Clozel--Delorme for the linear invariant Paley--Wiener theorem.  Li's theorem alone does not compare the two groups under study.
\item \textbf{Referee \S1.8: Proposition 9.1 only repeated the lattice criterion.} Merge it into a remark applying Proposition 6.1 (Part III, Corollary 5.2).  Do not represent it as a new conditional theorem.
\item \textbf{Referee \S2: Status labels, authors, citations, and versions inconsistent.} Adopt Parts I--III's [P]/[E]/[O]/[C] convention and update Part III's citation to Draft 9 with the actual full title and authors. Part I and II bibliography entries reflect the GitHub Draft 7 versions that were available at the time.
\item \textbf{Referee \S3: Editorial drafting history belongs in a log.} Remove the draft-by-draft cross-checking paragraph, the self-reference to Part IV Draft 1, and the phrase ``in the numbering cited by the referee.'' Trim repetitive disclaimers and clarify the abstract.
\end{enumerate}

\section*{The repaired Haar modulus computation}
Write $A_E=(\prod_i E_i^\times)/F^\times$, $K_0=\ker(P_n)$,
and take a small torsion-free compact open subgroup $U\subset A_E$
that meets $K_0$ trivially.  Its principal-unit logarithm has
$\mathbb Z_p$-rank $(r-1)[F:\mathbb Q_p]$, so
\[
[U:U^n]=|n|_F^{-(r-1)},\qquad
P_n^{-1}(U^n)=UK_0.
\]
Thus the pushforward Haar density on $P_n(A_E)$ equals
\[
\frac{\operatorname{vol}(UK_0)}{\operatorname{vol}(U^n)}
=|K_0|[U:U^n]=n^{k-1}|n|_F^{-(r-1)}.
\]
This is a calculation for the particular local group $A_E$,
not a universal quotient-index identity.  Although this factor
numerically equals $[A_E:P_n(A_E)]/|K_0|$ here, that expression
is not generally a Haar Jacobian.

\section*{Precisely what has and has not advanced in Section 8}
Flicker--Kazhdan \S6 already imposes a regular orbital-integral
invariance when different ordinary roots yield the same central
translate of a star image.  Draft 3 proves that a pointwise
covered-to-ordinary orbital-integral identity \emph{requires}
this invariance.  It then isolates an explicit set of sufficient
hypotheses for a geometric trace comparison, including local
constants $\lambda_v$, a global product identity, rational
splittings, quotient Haar measures, and support eliminating
singular classes.  The implication is proved once these inputs
are granted, but the existence of matching data satisfying them
is still open.  In particular, a local root sum of Zou's type
cannot automatically be identified with one rational orbital
integral, and regular roots with singular powers require separate
germ or support control.

\section*{Bibliographic checks}
The updated Part III reference is to: D. Joyner, ChatGPT (OpenAI),
and Claude (Anthropic), \emph{A metaplectic Jacquet--Langlands transfer
for division algebras of prime degree, III: The division-algebra side
and its local character transform}, Draft 9 (November 7, 2026).
The erroneous, unused Takeda entry from Draft 2 (2017, pp. 415--449)
has been removed; the verified journal pages are 199--230,
\url{https://doi.org/10.2140/pjm.2017.290.199}.
A Clozel--Delorme citation has been added:
\emph{Le th\'eor\`eme de Paley--Wiener invariant pour les groupes de
Lie r\'eductifs}, Invent. Math. 77 (1984), 427--453.
Li's statement is supported by his 2014 article,
\url{https://doi.org/10.5802/aif.2915}.

\section*{Outstanding mathematical questions}
\begin{enumerate}[leftmargin=*]
\item Wild-place integrality of the matrix $(a_{\rho\tau})$ remains
unproved; the concrete $r=2,n=3$ example should be computed in a
twisted finite character algebra, with level control.
\item Split-place transfer requires test functions satisfying the
FK compatibility and the pointwise $\lambda_v$-identity with the
chosen measures; the present work provides a necessity criterion,
not an existence theorem.
\item The product identity for the $\lambda_v$ and the ramified
$b_{D,v}$, together with germ behavior and global spectral isolation,
must be established before an ordinary-versus-covered $D^\times$
trace formula can be asserted.
\item Part III Corollary 7.8 is a proved tame character statement only
in the explicitly available Flicker--Kazhdan range (notably
$(n,N_r)=1$ as a sufficient published hypothesis); it does not
settle wild places.
\end{enumerate}

\section*{Compilation and distribution}
Draft 3 and this log were compiled using \texttt{pdflatex};
PDF outputs were inspected. The accompanying archive carries local
Eastern Standard Time file-entry timestamps (UTC-05:00) and a
README listing source baselines and limitations.  No GitHub push is
performed by generating the archive.
\end{document}
